\chapter{Planificación y Diseño}\label{cap:planificacion}

\section{Ideas y Alcance}

Esta fase inicial fue determinante para establecer los cimientos de este proyecto. El
primer hito fue una reunión inicial de definición de alcance donde se determinó que,
para que el sistema fuera útil, debía ser capaz de manejar distintos juegos de cartas
sin reestructurar la base de datos para cada franquicia.

Durante la reunión, contemplamos diversas formas de orientar el proyecto sobre la
base de las tecnologías escogidas. Tras analizar cómo usar APIs externas para la
extracción de cartas, llegamos a la conclusión de que la creación de barajas
personalizadas era la idea con más sentido para implementar el CRUD completo que se
nos pedía. La visión era crear una base sólida que permitiera, en las siguientes
fases, mejorar la lógica interna de validación de reglas.

\section{Análisis del Flujo de Usuario y Casos de Uso}

Para que la aplicación sea operativa, se definió un flujo de usuario lineal y lógico
que maximiza la eficiencia en la gestión de colecciones. Este flujo se documentó bajo
los siguientes pasos:

\begin{itemize}
	\item \textbf{Identificación y Autenticación:} El usuario accede al sistema, lo
	que permite vincular sus barajas personalizadas a su perfil único almacenado en
	MongoDB.
	\item \textbf{Consulta y Descubrimiento:} El usuario utiliza el buscador. Este
	componente realiza peticiones a la API externa de tcgdex.dev, devolviendo el
	listado de cartas candidatas.
	\item \textbf{Persistencia en Colección:} Una vez seleccionada una carta, el
	sistema CRUD realiza una operación de creación (CREATE) para guardar la referencia
	y los metadatos específicos.
	\item \textbf{Construcción de Barajas:} El usuario crea un contenedor
	``Baraja'' donde añade y retira cartas. Aquí se aplica la lógica de edición
	(UPDATE) para modificar cantidades o estados.
	\item \textbf{Ver y Eliminar:} El usuario puede ver sus barajas, lo que sería la
	funcionalidad de (READ), y podría, si quisiera, eliminarlas con la operación de
	(DELETE).
\end{itemize}

\section{Diseño de la Arquitectura de Datos en MongoDB}

Al haber seleccionado MongoDB por el interés del equipo en las bases de datos no
relacionales, el diseño de la arquitectura fue un punto clave. A diferencia de un
modelo SQL rígido, optamos por un esquema basado en documentos BSON, lo que permite
una flexibilidad total en los atributos.

\textbf{Figura 1.} Modelo datos User (MongoDB)

\textbf{Figura 2.} Modelo datos Deck (MongoDB)

\textbf{Figura 3.} Modelo datos User\_Card (MongoDB)

\section{Stack Tecnológico}

El requisito principal del proyecto era construir una API REST que permitiese trabajar
con una base de datos NoSQL en MongoDB. Para este proyecto en concreto, teníamos la
restricción de que la tecnología a utilizar debía ser C\#, por lo que construimos el
backend en .NET 8. Este se encarga de la lógica de negocio y exponer la API.

Por otro lado, para ampliar el alcance inicial, se decidió incluir una capa frontend
mediante React + Vite + TypeScript, con el objetivo de facilitar el uso de la API de
cara a los usuarios.

En cuanto a infraestructura, el proyecto se desplegó en servicios cloud:

\begin{itemize}
	\item Render para backend y frontend.
	\item MongoDB Atlas para alojar la base de datos en entorno gestionado.
\end{itemize}

De esta forma, la solución mantiene el núcleo requerido (MongoDB + C\#) y añade una
arquitectura web completa, preparada para uso real y validación en producción.

\section{Resultado de la fase}

La fase de planificación y diseño ha permitido establecer una base sólida para el
desarrollo del proyecto. Durante esta etapa se han definido claramente los objetivos
del sistema, el alcance funcional y el flujo de usuario, asegurando que la aplicación
cubra las necesidades principales de gestión de barajas.

Como resultado de esta fase, se ha obtenido una visión clara del funcionamiento global
del sistema, incluyendo la interacción entre usuario, aplicación y fuentes externas de
datos. Además, ha quedado definido el stack tecnológico. Esto ha facilitado la
transición hacia la fase de desarrollo, donde se ha podido implementar el sistema de
manera más eficiente y con menor riesgo de errores estructurales.
